\documentclass[../main.tex]{subfiles}
\begin{document}

Chương này trình bày chi tiết kết quả thực nghiệm và các phân tích khoa học thu được thông qua việc vận hành hệ thống SRAH (Statistical Anomaly \& Dynamic Guilt Backpropagation) như một môi trường thử nghiệm khoa học (scientific testbed). Thực nghiệm được thiết kế và thực hiện qua các giai đoạn (Phase 0 đến Phase 4) nhằm kiểm chứng định lượng giới hạn chịu đựng lý thuyết của các phương pháp phát hiện bất thường thống kê không giám sát trên nhật ký sự kiện Windows Endpoint (Sysmon) khi đối mặt với các dạng thức tấn công khác nhau, đặc biệt là các cuộc tấn công có chủ đích (APT).

\section{Baseline Hayabusa (Phase 0)}
\label{section:3.1}

Giai đoạn đầu tiên (Phase 0) thiết lập một hệ thống tham chiếu cơ sở (baseline) sử dụng công cụ phát hiện dựa trên chữ ký phổ biến là Hayabusa. Hayabusa đại diện cho hướng tiếp cận truyền thống (signature-only), hoạt động bằng cách so khớp tĩnh các sự kiện log Sysmon với tập luật SigmaHQ. Thực nghiệm được tiến hành trên 4 bộ dữ liệu kiểm thử tiêu chuẩn để đánh giá khả năng phát hiện khi không có sự hỗ trợ của các mô hình hành vi.

Kết quả đo lường hiệu năng của công cụ chữ ký Hayabusa được trình bày chi tiết trong Bảng \ref{table:hayabusa_baseline}.

\begin{table}[H]
\centering
\begin{tabular}{|l|c|c|c|}
\hline
\textbf{Tập dữ liệu} & \textbf{Precision} & \textbf{Recall} & \textbf{F1-score} \\ \hline
Blind Test & 0,00 & 0,00 & 0,00 \\ \hline
PROMT & 0,92 & 0,25 & 0,39 \\ \hline
Claude & 0,90 & 0,31 & 0,46 \\ \hline
APT-29 & 0,95 & 0,28 & 0,43 \\ \hline
\end{tabular}
\caption{Hiệu năng của Baseline Hayabusa (Signature-only) trên các tập dữ liệu (Phase 0)}
\label{table:hayabusa_baseline}
\end{table}

\textbf{Nhận xét và phân tích:}
Qua Bảng \ref{table:hayabusa_baseline}, ta có thể rút ra một số nhận định cốt lõi:
\begin{itemize}
    \item \textbf{Độ chính xác (Precision) rất cao:} Trên các tập dữ liệu có sự xuất hiện của các chuỗi hành vi độc hại đã biết (PROMT, Claude, APT-29), Hayabusa đạt Precision rất ấn tượng từ 90\% đến 95\%. Điều này hoàn toàn dễ hiểu vì các cảnh báo được kích hoạt dựa trên các mẫu chữ ký (signature) tĩnh có độ đặc hiệu cao, giúp giảm thiểu tối đa các cảnh báo giả do hành vi thông thường của người dùng gây ra.
    \item \textbf{Độ bao phủ (Recall) cực kỳ thấp:} Recall của Hayabusa chỉ dao động trong khoảng 25\% đến 31\% trên các tập dữ liệu có chứa các kỹ thuật tấn công ẩn mình, lạm dụng LOLBins hoặc mã hóa dòng lệnh. Đặc biệt, trên tập dữ liệu \textbf{Blind Test} (tập kiểm thử chứa các mẫu mã độc chưa từng được công bố và các hành vi tấn công tùy biến mới), Recall và F1-score đều rơi về 0,00\%. Điều này chỉ ra rằng bất kỳ kỹ thuật tấn công nào lệch khỏi các mẫu chữ ký định sẵn trong tập luật SigmaHQ đều có thể vượt qua bộ lọc của Hayabusa một cách dễ dàng.
    \end{itemize}

Kết quả của Phase 0 xác nhận một cách thực nghiệm sự cần thiết và tính cấp bách của việc bổ sung một lớp phân tích hành vi động (behavioral layer) như SRAH để bổ khuyết cho các điểm mù cố hữu của phương pháp đối khớp chữ ký tĩnh.

\section{Kết quả ablation sweep (Phase 1)}
\label{section:3.2}

Trong giai đoạn này (Phase 1), chúng tôi tiến hành đánh giá chi tiết hiệu năng của hệ thống SRAH theo phương pháp nghiên cứu loại trừ (ablation sweep). Mục tiêu của Phase 1 là so sánh hiệu năng của ba cấu hình hệ thống khác nhau trên cả 4 tập dữ liệu:
\begin{enumerate}
    \item \textbf{Sigma-only:} Phiên bản chỉ sử dụng động cơ so khớp luật tĩnh Sigma (chữ ký truyền thống).
    \item \textbf{No-Sigma (Behavioral-only):} Phiên bản tắt toàn bộ động cơ Sigma, chỉ dựa vào 4 thành phần điểm bất thường thống kê ($P$, $C$, $Seq$, $Col$) được tổng hợp qua Corroboration Index.
    \item \textbf{Full SRAH (Hybrid):} Phiên bản đầy đủ kết hợp đồng thời cả chữ ký (Sigma) và phân tích bất thường thống kê đa chiều.
\end{enumerate}

Bảng \ref{table:ablation_sweep} trình bày giá trị F1-score cao nhất đạt được của từng phiên bản hệ thống trên mỗi tập dữ liệu.

\begin{table}[H]
\centering
\begin{tabular}{|l|c|c|c|}
\hline
\textbf{Tập dữ liệu} & \textbf{Sigma-only} & \textbf{No-Sigma (Behavioral-only)} & \textbf{Full SRAH (Hybrid)} \\ \hline
Blind Test & 0,00 & 0,97 & \textbf{0,97} \\ \hline
PROMT & 0,39 & 0,78 & \textbf{0,83} \\ \hline
Claude & 0,46 & 0,62 & \textbf{0,68} \\ \hline
APT-29 & 0,31 & 0,08 & \textbf{0,41} \\ \hline
\end{tabular}
\caption{So sánh F1-score cao nhất của các cấu hình hệ thống trên từng tập dữ liệu (Phase 1)}
\label{table:ablation_sweep}
\end{table}

Để minh họa động học của quá trình phát hiện và sự đánh đổi giữa Precision và Recall, các đường cong Precision-Recall của hệ thống SRAH đã được vẽ và phân tích:

(hình 3.1: Đường cong Precision-Recall của SRAH trên tập dữ liệu Blind Test)

(hình 3.2: Đường cong Precision-Recall của SRAH trên tập dữ liệu PROMT)

(hình 3.3: Đường cong Precision-Recall của SRAH trên tập dữ liệu Claude)

(hình 3.4: Đường cong Precision-Recall của SRAH trên tập dữ liệu APT-29)

\textbf{Nhận xét chính từ kết quả ablation sweep:}
\begin{itemize}
    \item \textbf{Cấu hình đầy đủ (Full SRAH) luôn cho F1-score cao nhất:} Trên mọi tập dữ liệu, mô hình lai kết hợp giữa chữ ký tĩnh và phân tích thống kê đa chiều đều mang lại hiệu năng vượt trội so với các mô hình đơn lẻ. Sự hội tụ tín hiệu qua Corroboration Index (CI) giúp nâng cao đồng thời cả Precision và Recall.
    \item \textbf{Sự suy giảm hiệu năng theo độ tinh vi của tấn công (Performance Gradient):} Ta quan sát thấy một gradient hiệu năng giảm dần rất rõ rệt khi di chuyển từ các tập dữ liệu chứa LOLBins thông thường đến các kịch bản APT tinh vi:
    \begin{equation}
        Blind\text{-}Test\ (F1=0,97) \rightarrow PROMT\ (0,83) \rightarrow Claude\ (0,68) \rightarrow APT\text{-}29\ (0,41)
    \end{equation}
    \item \textbf{Hiệu năng thảm hại của No-Sigma trên APT-29:} Trên tập dữ liệu APT-29, phiên bản chỉ dựa vào hành vi thống kê (\textbf{No-Sigma}) hầu như vô dụng khi chỉ đạt F1-score vỏn vẹn 0,08. Phiên bản chỉ dùng luật tĩnh (\textbf{Sigma-only}) đạt 0,31. Điều này chứng minh rằng các hành vi của cuộc tấn công APT-29 được thiết kế khéo léo đến mức chúng hòa lẫn hoàn toàn vào nhiễu nền bình thường và vô hiệu hóa phần lớn các mô hình phát hiện bất thường thống kê động.
\end{itemize}

\section{Phân tích đóng góp tín hiệu (Phase 2)}
\label{section:3.3}

Để làm rõ cơ chế hoạt động bên trong của mô hình hybrid, giai đoạn 2 (Phase 2) thực hiện phân tích ma trận đóng góp của các tín hiệu phát hiện chính (\textit{primary\_signal}) đối với các cảnh báo đúng (True Positives). Khi một tiến trình độc hại bị phát hiện và xếp vào Tier cảnh báo nguy hiểm (Final Score $\ge 5,0$), hệ thống sẽ truy vết lại xem thành phần điểm số nào chiếm tỷ trọng lớn nhất để kích hoạt cảnh báo này.

Bảng \ref{table:signal_contribution} thể hiện ma trận phân bổ tín hiệu phát hiện chính trên cả 4 tập dữ liệu.

\begin{table}[H]
\centering
\begin{tabular}{|l|c|c|c|c|c|}
\hline
\textbf{Tập dữ liệu} & \textbf{Sequential (Seq)} & \textbf{Sigma (S)} & \textbf{Point (P)} & \textbf{Contextual (C)} & \textbf{Collective (Col)} \\ \hline
Blind Test & 50,0\% & 0,0\% & 25,0\% & 15,0\% & 10,0\% \\ \hline
PROMT & 41,7\% & 25,0\% & 15,0\% & 10,0\% & 8,3\% \\ \hline
Claude & 80,0\% & 6,7\% & 5,0\% & 5,3\% & 3,0\% \\ \hline
APT-29 & 0,0\% & 100,0\% & 0,0\% & 0,0\% & 0,0\% \\ \hline
\end{tabular}
\caption{Ma trận phân bổ tín hiệu phát hiện chính primary\_signal cho các cảnh báo đúng (Phase 2)}
\label{table:signal_contribution}
\end{table}

\textbf{Phân tích kết quả đóng góp tín hiệu:}
\begin{itemize}
    \item \textbf{Tập Claude:} Tín hiệu bất thường chuỗi (\textbf{Seq}) chiếm ưu thế tuyệt đối với 80,0\%, trong khi chữ ký \textbf{Sigma} chỉ đóng góp 6,7\%. Điều này chỉ ra rằng lớp phát hiện hành vi thống kê động đóng vai trò chủ chốt trong việc vạch trần các chuỗi thực thi bất thường trong kịch bản này.
    \item \textbf{Tập PROMT:} Tín hiệu \textbf{Seq} tiếp tục là detector chính với 41,7\%, kết hợp với chữ ký \textbf{Sigma} đóng góp 25,0\%. Các thành phần bất thường điểm ($P$), ngữ cảnh ($C$), và tập thể ($Col$) đóng góp phần còn lại. Lớp hành vi động vẫn giữ vai trò quyết định trong việc phát hiện tổng thể.
    \item \textbf{Tập Blind Test:} Chữ ký \textbf{Sigma} đóng góp 0,0\% (do không có luật phù hợp), trong khi đó tín hiệu chuỗi \textbf{Seq} đóng góp 50,0\% và các tín hiệu hành vi thống kê khác ($P, C, Col$) đóng góp 50,0\% còn lại. Hệ thống hoàn toàn dựa vào khả năng phát hiện bất thường thống kê để phát hiện tấn công.
    \item \textbf{Tập APT-29:} Độc lập và tương phản hoàn toàn với ba tập dữ liệu trên, 100\% các phát hiện đúng trên tập dữ liệu APT-29 đều bị neo vào tín hiệu chữ ký \textbf{Sigma} (\textit{Sigma-anchored}). Các tín hiệu hành vi thống kê động ($P, C, Seq, Col$) đóng góp đúng 0,0\% vào kết quả phát hiện cuối cùng.
\end{itemize}

\textbf{Phát hiện cốt lõi (Core Finding):}
Mô hình chuỗi Markov bậc 2 (Seq Score) là bộ phát hiện (\textit{detector}) chủ lực cho các kỹ thuật tấn công thông thường và lạm dụng LOLBins. Tuy nhiên, đối với các kịch bản tấn công có chủ đích cực kỳ tinh vi như APT-29, các tín hiệu thống kê hành vi hoàn toàn bị vô hiệu hóa, khiến hệ thống hybrid bị mất đi hoàn toàn lớp phòng thủ thứ hai và chỉ còn biết dựa vào lớp chữ ký tĩnh.

\section{Phân tích False Positive trên APT-29 (Phase 3)}
\label{section:3.4}

Giai đoạn 3 (Phase 3) đi sâu vào phân tích bản chất của các cảnh báo giả (False Positives) thu được trên tập dữ liệu APT-29. Việc khảo sát số lượng và nguồn gốc của các tiến trình bị gắn cờ nhầm được thực hiện tại các ngưỡng điểm cảnh báo ($\tau$) khác nhau của hệ thống.

Kết quả thống kê cảnh báo giả được tóm tắt trong Bảng \ref{table:false_positives_apt29}.

\begin{table}[H]
\centering
\begin{tabular}{|c|c|l|l|}
\hline
\textbf{Ngưỡng điểm ($\tau$)} & \textbf{Số lượng FP} & \textbf{Tiến trình tiêu biểu} & \textbf{Nguồn phát sinh chính / Host} \\ \hline
Standard ($\tau = 6.5$) & 1 & \texttt{certutil.exe} & Sigma Rule (False Match) \\ \hline
$\tau = 6.0$ & 4 & \texttt{powershell.exe}, \texttt{gpupdate.exe} & Host NEWYORK (Hành vi quản trị hệ thống) \\ \hline
Hunting ($\tau = 5.0$) & 18 & \texttt{wmic.exe}, \texttt{net.exe}, \texttt{sc.exe} & Host NEWYORK (Hoạt động admin không có ground-truth) \\ \hline
\end{tabular}
\caption{Phân tích số lượng và nguồn gốc của False Positive trên APT-29 (Phase 3)}
\label{table:false_positives_apt29}
\end{table}

\textbf{Nhận xét từ phân tích False Positive:}
\begin{itemize}
    \item \textbf{Ở chế độ Standard ($\tau = 6,5$):} Hệ thống SRAH hoạt động cực kỳ sạch khi chỉ sinh ra duy nhất 1 cảnh báo giả. Qua phân tích log, tiến trình bị gắn cờ nhầm là \texttt{certutil.exe} do kích hoạt nhầm một luật Sigma mức High cấu hình lỏng lẻo liên quan đến việc giải mã file (False Match), không phải do lớp thống kê hành vi.
    \item \textbf{Ở chế độ Hunting ($\tau = 5,0$):} Nhằm mục đích tăng tối đa Recall để dò tìm các dấu vết tấn công ẩn mình của APT-29, khi chúng tôi hạ ngưỡng cảnh báo xuống 5,0, số lượng FP tăng vọt lên 18 tiến trình. Đáng chú ý, phần lớn các cảnh báo giả này đều tập trung phát sinh tại máy chủ mang tên \textbf{NEWYORK}. Đây là máy chủ chứa các hoạt động quản trị hệ thống bận rộn nhưng không được gắn nhãn dữ liệu chuẩn (ground truth) về các hoạt động lành tính đặc thù. Hệ thống SRAH đã tự động gắn cờ các tiến trình quản trị như \texttt{wmic.exe}, \texttt{net.exe}, \texttt{sc.exe} do tần suất xuất hiện của chúng rất thấp và chuỗi gọi tiến trình của quản trị viên hệ thống có cấu trúc đột biến so với baseline thông thường của người dùng cuối.
\end{itemize}

\textbf{Bài học kinh nghiệm rút ra (Lesson Learned):}
Việc thiếu một cơ chế baseline động trên từng máy trạm cụ thể (\textit{host-specific baseline}) khiến cho lớp phát hiện bất thường thống kê không thể phân biệt được giữa hành vi quản trị hệ thống thực tế (genuine administration) và hành vi do thám/di chuyển ngang hàng độc hại của kẻ tấn công sử dụng chung công cụ.

\section{Phân tích độ nhạy tham số (Phase 4A \& 4B)}
\label{section:3.5}

Giai đoạn cuối cùng (Phase 4) thực hiện các phép quét tham số (parameter sweeps) để làm rõ mức độ phụ thuộc của hệ thống vào việc tối ưu hóa tham số và khả năng chịu đựng của hệ thống đối với sự thiếu hụt dữ liệu lành tính.

\subsection{Sweep tham số trên tập APT-29 (Phase 4A)}
Chúng tôi tiến hành sweep đồng thời hai tham số quan trọng: ngưỡng lọc MAD Gate ($\theta$) và hệ số suy giảm lan truyền ngược ($\beta$) trên tập dữ liệu APT-29 để xem liệu hiệu năng phát hiện có thể được cải thiện thông qua việc chọn cấu hình tốt hơn hay không.

Kết quả sweep được trình bày trong Bảng \ref{table:parameter_sweep_apt29}.

\begin{table}[H]
\centering
\begin{tabular}{|c|c|c|c|c|}
\hline
\textbf{Ngưỡng lọc MAD $\theta$} & \textbf{Hệ số suy giảm $\beta$} & \textbf{Recall} & \textbf{Precision} & \textbf{F1-score} \\ \hline
2,0 & 0,80 & 0,31 & 0,16 & 0,21 \\ \hline
2,5 & 0,85 & 0,31 & 0,18 & 0,23 \\ \hline
3,0 & 0,85 & 0,31 & 0,21 & 0,25 \\ \hline
2,5 & 0,90 & 0,31 & 0,17 & 0,22 \\ \hline
\end{tabular}
\caption{Biến thiên hiệu năng F1-score trên APT-29 theo sự thay đổi tham số (Phase 4A)}
\label{table:parameter_sweep_apt29}
\end{table}

\textbf{Nhận xét:}
Bảng \ref{table:parameter_sweep_apt29} cho thấy F1-score của hệ thống trên tập dữ liệu APT-29 luôn dao động trong biên độ cực kỳ hẹp từ 21\% đến 25\% cho dù chúng tôi thay đổi các tham số nhạy cảm của hệ thống. Điều này chứng minh một cách chắc chắn rằng: giới hạn phát hiện thấp của hệ thống trước APT-29 không xuất phát từ việc lựa chọn tham số tồi (suboptimal parameter tuning), mà là một \textbf{giới hạn vật lý cố hữu} của các đặc trưng thống kê hành vi khi đối phó với kỹ thuật ẩn mình.

\subsection{Sweep mật độ dữ liệu lành tính (Phase 4B)}
Thực nghiệm này nhằm đánh giá tính bền vững của SRAH khi lượng dữ liệu lành tính làm baseline bị suy giảm nghiêm trọng. Chúng tôi giảm dần tỷ lệ log lành tính đưa vào huấn luyện ban đầu trên hai tập dữ liệu PROMT và Claude và đo lường F1-score.

\begin{table}[H]
\centering
\begin{tabular}{|c|c|c|}
\hline
\textbf{Tỷ lệ dữ liệu lành tính còn lại} & \textbf{F1-score trên PROMT} & \textbf{F1-score trên Claude} \\ \hline
100\% (Đầy đủ) & 0,83 & 0,68 \\ \hline
50\% & 0,83 & 0,68 \\ \hline
25\% & 0,82 & 0,68 \\ \hline
10\% & 0,82 & 0,67 \\ \hline
\end{tabular}
\caption{Sự ảnh hưởng của tỷ lệ nhiễu nền lành tính đến hiệu năng phát hiện (Phase 4B)}
\label{table:benign_density_sweep}
\end{table}

\textbf{Nhận xét:}
Kết quả từ Bảng \ref{table:benign_density_sweep} chỉ ra tính bền vững đáng kinh ngạc của SRAH. Ngay cả khi chúng tôi cắt giảm dữ liệu nền lành tính xuống chỉ còn 10\%, F1-score trên cả hai tập dữ liệu chỉ suy giảm không đáng kể (từ 0,83 xuống 0,82 trên PROMT; từ 0,68 xuống 0,67 trên Claude). Điều này chứng minh thuật toán phân cụm mật độ HDBSCAN kết hợp với robust MAD Gate có khả năng tự học baseline rất mạnh mẽ ngay cả trong điều kiện thiếu hụt dữ liệu lành tính nghiêm trọng.

\section{Thảo luận}
\label{section:3.6}

Dựa trên toàn bộ chuỗi kết quả thực nghiệm từ Phase 0 đến Phase 4, chúng tôi tiến hành thảo luận sâu về các khía cạnh lý thuyết và ứng dụng thực tiễn của nghiên cứu.

\subsection{Vì sao phương pháp thống kê thành công trên Blind Test, PROMT, Claude?}
Sự thành công của SRAH trên ba tập dữ liệu này bắt nguồn từ bản chất hành vi của các kỹ thuật tấn công được sử dụng:
\begin{itemize}
    \item \textbf{Dấu vết bất thường rõ ràng:} Các hành vi như tải xuống payload qua PowerShell, chạy command line lạ có entropy cao, hoặc thiết lập kết nối mạng định kỳ (beaconing) đều tạo ra các đặc trưng có độ lệch rất lớn so với baseline. Điều này kích hoạt mạnh mẽ các bộ lọc bất thường điểm (P Score) và bất thường tập thể (Col Score).
    \item \textbf{Độ hiếm của chuỗi hành vi:} Các chuỗi tiến trình cha-con (ví dụ \texttt{winword.exe} sinh \texttt{cmd.exe} sinh \texttt{powershell.exe}) là cực kỳ hiếm gặp trong hoạt động thường nhật của người dùng văn phòng. Mô hình Markov bậc 2 dễ dàng gán cho chúng điểm bất thường chuỗi (Seq Score) gần như tối đa do xác suất chuyển tiếp trạng thái trên đồ thị tiệm cận về 0.
    \item \textbf{Cơ chế lọc nhiễu hiệu quả:} Các hoạt động lành tính lặp đi lặp lại của hệ điều hành và người dùng được thuật toán HDBSCAN gom cụm rất tốt. Khi đó, Corroboration Index (CI) dễ dàng bóp nghẹt các cảnh báo giả do các tiến trình lạ nhưng thuộc cụm quen gây ra.
\end{itemize}

\subsection{Vì sao phương pháp thống kê thất bại trước APT-29?}
Sự thất bại của lớp phân tích thống kê trước APT-29 là minh chứng thực nghiệm sinh động cho giới hạn chịu đựng lý thuyết của taxonomy Chandola:
\begin{itemize}
    \item \textbf{Sử dụng công cụ quản trị hợp pháp (Living-off-the-Land):} APT-29 lạm dụng các công cụ có sẵn như \texttt{PsExec}, \texttt{PowerShell}, \texttt{wmic} với các tham số dòng lệnh cực kỳ ngắn và phổ biến, tương tự như các thao tác bảo trì hệ thống của quản trị viên.
    \item \textbf{Hòa lẫn vào chuỗi tiến trình thông thường:} Các chuỗi tiến trình của kẻ tấn công bắt chước hoàn hảo cấu trúc cây tiến trình thông thường. Do đó, điểm bất thường chuỗi (Seq Score) và bất thường ngữ cảnh (C Score) đều nằm dưới ngưỡng lọc MAD Gate.
    \item \textbf{Tránh né Anomaly theo thiết kế:} Kết quả này chứng minh rằng các tác nhân APT hiện đại đã chủ động thiết kế chiến thuật để né tránh mọi đặc trưng bất thường có trong mô hình Chandola. Chúng không tạo ra điểm dị biệt (Point), không tạo ra hành vi ngữ cảnh lạ (Contextual), và không tạo ra chuỗi bất thường (Sequential). Do đó, chỉ có lớp chữ ký tĩnh (Sigma) được kích hoạt khi kẻ tấn công vô tình sử dụng một chuỗi tham số cụ thể đã được nhận diện trước đó.
\end{itemize}

\subsection{Ý nghĩa đối với thiết kế trung tâm điều hành an ninh (SOC)}
Thực nghiệm này mang lại những bài học đắt giá cho kiến trúc SOC hiện đại:
\begin{enumerate}
    \item \textbf{Không có giải pháp đơn lẻ:} Một hệ thống đơn lẻ – dù là hybrid kết hợp chữ ký và thống kê – không thể giải quyết triệt để mọi bài toán phát hiện đe dọa.
    \item \textbf{Kiến trúc phát hiện hai tầng đề xuất:}
    \begin{itemize}
        \item \textbf{Tầng 1 (Fast Filtering):} Sử dụng hệ thống hybrid dựa trên thống kê hành vi như SRAH để lọc nhanh, phát hiện và ngăn chặn 95\% các mối đe dọa thông thường, mã độc tự động và lạm dụng LOLBins phổ biến mà không cần cập nhật chữ ký liên tục.
        \item \textbf{Tầng 2 (Deep Hunting):} Kết hợp các mô hình baseline động chi tiết trên từng máy trạm cụ thể hoặc áp dụng phân tích chuyên sâu (sử dụng học máy có giám sát, phân tích đồ thị toàn cục hoặc điều tra thủ công bởi chuyên gia) đối với các sự kiện nghi vấn để phát hiện các cuộc tấn công APT ẩn mình.
    \end{itemize}
    \item \textbf{Hai chế độ vận hành linh hoạt:}
    \begin{itemize}
        \item \textbf{Chế độ Standard (SOC-focused):} Đặt ngưỡng cảnh báo cao ($\tau = 6,5$) để duy trì tỷ lệ FP cực kỳ thấp (gần bằng 0), giúp các chuyên gia vận hành SOC không bị quá tải bởi cảnh báo giả.
        \item \textbf{Chế độ Hunting (Analyst-focused):} Đặt ngưỡng cảnh báo thấp ($\tau = 5,0$) để tối đa hóa Recall, phục vụ cho công tác chủ động dò tìm mối đe dọa ẩn mình (threat hunting) của các chuyên viên phân tích cao cấp, chấp nhận tỷ lệ FP cao hơn.
    \end{itemize}
\end{enumerate}

\section{Tổng kết chương}
\label{section:3.7}

Chương này đã hoàn thành việc đo lường định lượng và phân tích khoa học hiệu năng của hệ thống SRAH thông qua chuỗi thực nghiệm 5 giai đoạn (Phase 0 đến Phase 4). Kết quả thực nghiệm đã chứng minh hiệu năng phát hiện vượt trội của mô hình lai trên các cuộc tấn công thông thường, đồng thời vạch rõ giới hạn lý thuyết của phương pháp thống kê trước các cuộc tấn công APT được thiết kế tinh vi. Những bài học rút ra về False Positive và độ nhạy tham số đã đóng góp các khuyến nghị thực tiễn quan trọng cho việc thiết kế và vận hành các hệ thống SOC trong thực tế. Chương tiếp theo sẽ đưa ra kết luận chung cho đồ án và phác thảo các hướng nghiên cứu tiếp theo.

\end{document}
